% from ega0-1-fr.tex
\oldpage[0\textsubscript{I}]{11}
\begin{env}[1.0.1]
\label{0.1.0.1-fr}
Tous les anneaux considérés dans ce Traité posséderont un \emph{élément unité};
tous les modules sur un tel anneau seront supposés \emph{unitaires}; les homomorphismes d'anneaux seront toujours supposés \emph{transformer l'élément unité en élément unité}; sauf mention expresse du contraire, un sous-anneau d'un anneau $A$ sera supposé \emph{contenir l'élément unité de $A$}.
Nous considérerons surtout des anneaux \emph{commutatifs}, et lorsque nous parlerons d'anneau sans préciser, il sera sous-entendu qu'il s'agit d'un anneau commutatif.
Si $A$ est un anneau non nécessairement commutatif, par $A$-module nous entendrons toujours un module \emph{à gauche}, sauf mention expresse du contraire.
\end{env}

\begin{env}[1.0.2]
\label{0.1.0.2-fr}
Soient $A$, $B$ deux anneaux non nécessairement commutatifs, $\varphi:A\to B$ un homomorphisme.
Tout $B$-module à gauche (resp. à droite) $M$ peut être muni d'une structure de $A$-module à gauche (resp. à droite) en posant $a.m=\varphi(a).m$ (resp. $m.a=m.\varphi(a)$); lorsqu'il sera nécessaire de distinguer sur $M$ les structures de $A$-module et de $B$-module, nous désignerons par $M_{[\varphi]}$ le $A$-module à gauche (resp. à droite) ainsi défini.
Si $L$ est un $A$-module, un homomorphisme $u:L\to M_{[\varphi]}$ est donc un homomorphisme de groupes commutatifs tel que $u(a.x)=\varphi(a).u(x)$ pour $a\in A$, $x\in L$; on dira aussi que c'est un \emph{$\varphi$-homomorphisme} $L\to M$, et que le couple $(\varphi,u)$ (ou, par abus de langage, $u$) est un \emph{di-homomorphisme} de $(A,L)$ dans $(B,M)$.
Les couples $(A,L)$ formés d'un anneau $A$ et d'un $A$-module $L$ forment donc une \emph{catégorie} pour laquelle les morphismes sont les di-homomorphismes.
\end{env}

\begin{env}[1.0.3]
\label{0.1.0.3-fr}
Sous les hypothèses de (1.0.2), si $\mathfrak{J}$ est un idéal à gauche (resp. à droite) de $A$, nous noterons $B\mathfrak{J}$ (resp. $\mathfrak{J}B$) l'idéal à gauche (resp. à droite) $B\varphi(\mathfrak{J})$ (resp. $\varphi(\mathfrak{J})B$) de $B$ engendré par $\varphi(\mathfrak{J})$; c'est aussi l'image de l'homomorphisme canonique $B\otimes_A\mathfrak{J}\to B$ (resp. $\mathfrak{J}\otimes_A B\to B$) de $B$-modules à gauche (resp. à droite).
\end{env}

\begin{env}[1.0.4]
\label{0.1.0.4-fr}
Si $A$ est un anneau (commutatif), $B$ un anneau non nécessairement commutatif, la donnée d'une structure de \emph{$A$-algèbre} sur $B$ équivaut à la donnée d'un homomorphisme d'anneaux $\varphi:A\to B$ tel que $\varphi(A)$ soit contenu dans le centre de $B$.
Pour tout idéal $\mathfrak{J}$ de $A$, $\mathfrak{J}B=B\mathfrak{J}$ est alors un idéal bilatère de $B$, et pour tout $B$-module $M$, $\mathfrak{J}M$ est alors un $B$-module égal à $(B\mathfrak{J})M$.
\end{env}

\begin{env}[1.0.5]
\label{0.1.0.5-fr}
Nous ne reviendrons pas sur les notions de \emph{module de type fini} et d'\emph{algèbre} (commutative) \emph{de type fini}; dire qu'un $A$-module $M$ est de type fini signifie qu'il existe

